% ===== 原文 pp.51--56 =====

% [原文 p.51]

这些涡旋位形不能提升为 $\cG$ 的涡旋位形。群 $G_{\min}$ 是李代数为 $\fg$、且允许物质场 $\phi_i$ 的表示 $R_i$ 的极小连通李群。在插入 $U_g^{(e)}(\Sigma_{d-2})$ 的周围我们有
\begin{equation}
    \int_{D^2} w_2^{\min} = g \,,
    \tag{3.156}
\end{equation}
其中 $w_2^{\min}$ 是把 $G_{\min}$ 丛提升为 $\cG$ 丛的障碍类。

生成磁 $(d-3)$-形式对称性的拓扑面算符仍是 (3.134)。

\subsubsection{多个规范群因子}

我们现在可以容易地把上述论证推广，以确定具有一般非阿贝尔规范代数
\begin{equation}
    \fg = \bigoplus_a \fg_a
    \tag{3.157}
\end{equation}
的规范理论中的电与磁高阶形式对称性；其中每个 $\fg_a$ 都是单的，因而取值于 (3.114)。这类理论包括\textbf{箭图规范理论}（quiver gauge theory），它们在量子场论的多个研究领域中扮演重要角色。

\para{规范群。}与 $\fg$ 相伴的单连通群是
\begin{equation}
    G = \prod_a G_a \,,
    \tag{3.158}
\end{equation}
其中 $G_a$ 是与每个 $\fg_a$ 相伴的单连通群。

我们想要考虑规范群为连通李群 $\cG$、且其李代数为 $\fg$ 的情形
\begin{equation}
    \mathrm{Lie}(\cG) = \fg \,.
    \tag{3.159}
\end{equation}
和之前一样，这样的群可以表示为
\begin{equation}
    \cG = \frac{G}{\cZ} \,,
    \tag{3.160}
\end{equation}
其中 $\cZ$ 是 $G$ 的中心 $Z(G)$ 的子群
\begin{equation}
    \cZ \subseteq Z(G) = \prod_a Z(G_a) \,.
    \tag{3.161}
\end{equation}
规范群 $\cG$ 的中心 $Z(\cG)$ 可以像 (3.121) 那样表示。

% [Tutor Note: p.51 式 (3.157)--(3.161) 注意区分：规范代数是直和 ⊕（不同单因子的场不互相混合），而单连通群是直积 ∏；相应地中心也是直积，故每个因子各贡献一份离散中心 Z_{n_a}。箭图规范理论正是这一结构在物理上的实现：每个节点给一个规范群因子、每条边给一个双基础物质场。]

% [原文 p.52]

\para{物质场。}我们引入物质场（它们可以具有任意的时空变换性质与任意质量）
\begin{equation}
    \phi_i ,\qquad 1 \le i \le m \,,
    \tag{3.162}
\end{equation}
它们按规范群 $\cG$ 的不可约表示 $R_i$ 变换。我们可以把这些表示写成张量积
\begin{equation}
    R_i = \bigotimes_a R_{i,a} \,,
    \tag{3.163}
\end{equation}
其中
\begin{equation}
    R_{i,a} = \text{$G_a$ 的一个不可约表示} .
    \tag{3.164}
\end{equation}
然而，要使 $R_i$ 成为 $\cG$ 的表示，对 $R_{i,a}$ 还有约束：要做到这一点，子群 $\cZ$ 必须在 $R_i$ 上作用平凡。

\para{高阶形式对称性。}电与磁 1-形式对称性的分析与上面类似。最终，电 1-形式对称性可以表示为 (3.140) 或 (3.142)，而磁 $(d-3)$-形式对称性表示为 (3.130)。这些高阶形式对称性的拓扑算符与 3.3.4 节所述的相同。

(3.139) 的计算可以用 Smith 标准形（Smith Normal Form）容易地完成。首先，我们可以把 $\widehat{Z}(\cG)$ 表示为
\begin{equation}
    \widehat{Z}(\cG) = \prod_{\alpha=1}^{p} \Z_{n_\alpha} \,,
    \tag{3.165}
\end{equation}
于是每个荷可以表示为
\begin{equation}
    q(R_i) = (q_{1,i}, \cdots, q_{\alpha,i}, \cdots, q_{p,i}) ,\qquad q_{\alpha,i} \in \{0,1,\cdots,n_\alpha-1\} .
    \tag{3.166}
\end{equation}
由此定义一个 $p \times m$ 矩阵 $Q$，其 $(\alpha,i)$ 元为 $q_{\alpha,i}$；再用它定义一个 $2p \times (m+p)$ 矩阵
\begin{equation}
    Y = (Z|Q) \,,
    \tag{3.167}
\end{equation}
其中 $Z$ 是对角的 $p \times p$ 矩阵，其第 $\alpha$ 个对角元是 $n_\alpha$。于是我们有
\begin{equation}
    G^{(1)} = \prod_{i=1}^{\min(2p,m+p)} \widehat{\Z}_{q_i^{\mathsf{D}}} \,,
    \tag{3.168}
\end{equation}
其中 $q_i^{\mathsf{D}}$ 由
\begin{equation}
    \diag(q_1^{\mathsf{D}}, q_2^{\mathsf{D}}, \cdots, q_{\min(2p,m+p)}^{\mathsf{D}}) \subset \mathrm{SNF}(Y)
    \tag{3.169}
\end{equation}
定义。

% [Tutor Note: p.52 式 (3.167) 的矩阵尺寸与原文不合：Z 是 p×p、Q 是 p×m，故 (Z|Q) 只能是 p×(m+p)，不可能是原文所写的 2p×(m+p)；例 3.11 中 p = 2、m = 1，而 (3.174) 的 Y 恰为 2×3 = p×(m+p)。相应地 (3.168)--(3.169) 中的 min(2p,m+p) 疑应为 min(p,m+p) = p。此处按原文照译，并请读者留意这一疑似笔误。]

% [Tutor Note: p.52 式 (3.165)：有限 阿贝尔群的分解 \hat{Z}(G) = ∏ Z_{n_α} 是 (3.168) 里逐因子取对偶的前提；Smith 标准形对整数矩阵作行列变换使之对角化，其对角元 q_i^D 给出商群 \hat{Z}(G)/Q_matter 的不变量因子。]

% [原文 p.53]

\begin{examplebox}{例 3.11：一个正交辛箭图规范理论}
考虑一个 $d$ 维箭图规范理论，其规范群为
\begin{equation}
    \cG = \Spin(4n+2) \times \Sp(m) \,,
    \tag{3.170}
\end{equation}
物质场处于 $\cG$ 的双基础表示
\begin{equation}
    R = V \otimes F \,,
    \tag{3.171}
\end{equation}
其中 $V$ 是 $\Spin(4n+2)$ 的矢量表示，$F$ 是 $\Sp(m)$ 的基础表示。我们有
\begin{equation}
    \widehat{Z}(\cG) = \widehat{Z}(\Spin(4n+2)) \times \widehat{Z}(\Sp(m)) = \Z_4 \times \Z_2 \,,
    \tag{3.172}
\end{equation}
而 $R$ 的荷是
\begin{equation}
    q(R) = (2,1) \,,
    \tag{3.173}
\end{equation}
由此我们构造出矩阵
\begin{equation}
    Y = \begin{pmatrix} 4 & 0 & 2 \\ 0 & 2 & 1 \end{pmatrix} ,
    \tag{3.174}
\end{equation}
它的 Smith 标准形是
\begin{equation}
    \mathrm{SNF}(Y) = \begin{pmatrix} 0 & 4 & 0 \\ 0 & 0 & 1 \end{pmatrix} ,
    \tag{3.175}
\end{equation}
这意味着该箭图规范理论的电 1-形式对称性是
\begin{equation}
    G^{(1)} = \Z_4 \,.
    \tag{3.176}
\end{equation}
磁 $(d-3)$-形式对称性是平凡的，因为 $\cG$ 是单连通的。
\end{examplebox}

% [Tutor Note: p.53 例 3.11 的 SNF 可以直接手算复核：(3.174) 的矩阵元素的最大公约数是 1，故第一个不变因子 d1 = 1；所有 2×2 子式的最大公约数是 gcd(8,4,4) = 4，故 d2 = 4。于是不变量因子为 (1,4)，平凡因子 1 被丢掉，剩下 Z4，与 (3.176) 一致。原文把对角元画在错开的位置上（(0,4,0),(0,0,1)），只是排版约定。]

\subsection{四维规范理论中的离散 theta 角}

在 $d=4$ 维时空中，有更多种类的带离散 theta 角的非阿贝尔规范理论 [42]。它们一般可以具有双荷 1-形式对称性。本小节的目的就是描述这些 1-形式对称性。

\subsubsection{纯 $\SO(3)_-$ Yang--Mills 理论}

在上一小节中，我们在一般的 $d$ 中讨论了纯 $\SU(2)$ 与 $\SO(3)$ Yang--Mills 理论，二者都基于规范代数 $\fsu(2)$。在 $d=4$ 中，还存在另一种基于 $\fsu(2)$ 规范代数的纯 Yang--Mills 理论 [37]，它具有
\begin{equation}
    \cG = \SO(3)
    \tag{3.177}
\end{equation}

% [原文 p.54]

的规范群，被称为 $\SO(3)_-$ 理论 [42]。相应地，上文讨论的通常的 $\SO(3)$ 理论有时也被称为 $\SO(3)_+$ 理论。

\para{线算符。}让我们先回顾 $\SO(3)_+$ 理论中关于线算符的讨论。首先，我们有纯 Wilson 线算符
\begin{equation}
    \cW_j ,\qquad j \in \Z \,.
    \tag{3.178}
\end{equation}
此外，我们还有在包围它们的球面 $S^2$ 上诱导出 $\SO(3)$ 磁单极子位形的线算符。考虑由余特征标描述的磁单极子位形
\begin{equation}
    \phi:\ \U(1) \to \SO(3) \,.
    \tag{3.179}
\end{equation}
诱导该磁单极子位形 $\phi$ 的各种线算符由 $\fsu(2)$ 的整数自旋表示标记。我们把它们记作
\begin{equation}
    \cD_{j,\phi} ,\qquad j \in \Z \,.
    \tag{3.180}
\end{equation}
特别地，对于 $\fsu(2)$ 的平凡表示 $j=0$，我们有一个线算符，称之为纯 't Hooft 线算符，并记作
\begin{equation}
    \cH_\phi \equiv \cD_{0,\phi} \,.
    \tag{3.181}
\end{equation}
注意 $\cD_{j,\phi}$ 出现在 $\cW_j$ 与 $\cH_\phi$ 的 OPE 中，我们通过把二者等同来表达这一事实
\begin{equation}
    \cD_{j,\phi} \equiv \cW_j \cH_\phi \,.
    \tag{3.182}
\end{equation}

\begin{notebox}{OPE 的记号}
文献中常用加法记号表示 OPE，此时所得的线算符写作 $\cW + \cH$。本文采用式 (3.182) 那样的乘法记号，以避免融合积与线的直和之间的混淆。两条线 $L_1$ 与 $L_2$ 的融合积 $L_1 \otimes L_2$ 的关联函数是
\begin{equation}
    \langle L_1 L_2 \rangle \,,
    \tag{3.183}
\end{equation}
而直和 $L_1 \oplus L_2$ 的关联函数是
\begin{equation}
    \langle L_1 \rangle + \langle L_2 \rangle \,.
    \tag{3.184}
\end{equation}
\end{notebox}

对于 $\SO(3)_-$ 理论，我们有同一组 (3.178) 的纯 Wilson 线算符。然而，诱导磁单极子位形的线算符谱不同。如果 $\phi$

% [原文 p.55]

可以被提升为 $\SU(2)$ 的余特征标，即如果相应的 $\SO(3)$ 磁单极子位形可以提升为 $\SU(2)$ 磁单极子位形，那么双荷线与 $\SO(3)_+$ 理论中的相同
\begin{equation}
    \cD_{j,\phi} ,\qquad j \in \Z \,.
    \tag{3.185}
\end{equation}
特别地，对于这样的余特征标 $\phi$，我们仍然有纯 't Hooft 算符
\begin{equation}
    \cH_\phi \,,
    \tag{3.186}
\end{equation}
然而，如果 $\phi$ 不能被提升为 $\SU(2)$ 的余特征标，即如果相应的 $\SO(3)$ 磁单极子位形不能提升为 $\SU(2)$ 磁单极子位形，那么双荷算符必须带有 $\fsu(2)$ 的半整数自旋表示，我们把它们记作
\begin{equation}
    \cD_{j,\phi} ,\qquad j \notin \Z \,.
    \tag{3.187}
\end{equation}
特别地，对于这样的余特征标 $\phi$，我们不再有纯 't Hooft 算符。我们可以形式地把这样的线算符写成
\begin{equation}
    \cD_{j,\phi} \equiv \cW_j \cH_\phi ,\qquad j \notin \Z \,,
    \tag{3.188}
\end{equation}
但应注意，$\cW_j$ 与 $\cH_\phi$ 各自都不是真线算符\footnote{正如我们将在后文讨论的，半整数 $j$ 的 $\cW_j$ 与非 $\SU(2)$ 的 $\phi$ 的 $\cH_\phi$ 都作为非真线算符存在于 $\SO(3)_-$ 理论中。}。然而，它们的乘积 $\cD_{j,\phi}$ 是真线算符。

% [Tutor Note: p.54--55 的记号 W_j、D_{j,φ}、H_φ 都是花体：W_j 是纯 Wilson 线、H_φ 是纯 't Hooft 线、D_{j,φ} 是双荷线（'t Hooft--Wilson 线）。式 (3.181) 的 H_φ ≡ D_{0,φ} 与 (3.182) 的 D_{j,φ} ≡ W_j H_φ 合起来说明：先把纯 't Hooft 线定义为 j = 0 的双荷线，再把一般双荷线定义为 Wilson 线与 't Hooft 线的融合积。]

\para{Witten 效应。}理解具有上述线算符谱的 $\SO(3)_-$ 理论存在性的一种有趣方式，是从 $\SO(3)_+$ 理论出发，把通常的（相对于离散的）theta 角改变为
\begin{equation}
    \theta \to \theta + 2\pi \,.
    \tag{3.189}
\end{equation}
众所周知，对 theta 角的这种变换会诱导 Witten 效应：按此效应，带磁荷的线算符会获得额外的电荷；当 $\phi$ 不能被提升为 $\SU(2)$ 余特征标时，这恰好把 $\SO(3)_+$ 理论的线算符 (3.182) 变成线算符 (3.188)。由于线算符 (3.188) 作为真算符在 $\SU(2)$ 与 $\SO(3)_+$ 理论中并不存在，我们得出结论：我们已经来到了一个新理论，我们称它为 $\SO(3)_-$ 理论。

% [Tutor Note: p.55 的 Witten 效应论证可用半经典公式核对：磁荷 m、电荷 e 的线算符在 θ → θ + 2π 下获得 e → e + m（因为作用量里 θ 项为 (θ/2π)∫F，量子化要求 ∫F/2π ∈ Z）。SO(3) 中不能提升到 SU(2) 的余特征标对应 m 为奇数，这正是半整数 j 的 (3.188) 出现的原因。]

% [原文 p.56]

\para{屏蔽。}在 $\SO(3)_-$ 理论中线算符的屏蔽与 $\SO(3)_+$ 理论中的屏蔽类似。所有纯 Wilson 算符与纯 't Hooft 算符都被完全屏蔽。另一方面，'t Hooft--Wilson 算符 (3.188) 也全都由屏蔽相互关联
\begin{equation}
    \cD_{j,\phi} \sim \cD_{j',\phi'} \,,
    \tag{3.190}
\end{equation}
但并未被完全屏蔽。因此 $\SO(3)_-$ 理论有一个由未被屏蔽线算符构成的群
\begin{equation}
    \mathsf{D}_1 = \Z_2 \,,
    \tag{3.191}
\end{equation}
这意味着该理论有一个 1-形式对称性群
\begin{equation}
    G^{(1)} = \widehat{\mathsf{D}}_1 = \Z_2 \,.
    \tag{3.192}
\end{equation}
't Hooft--Wilson 线算符 (3.188) 在 $G^{(1)}$ 下带非平凡的荷，而纯 Wilson 算符与纯 't Hooft 算符在它下面不带荷。因此，我们可以把 $G^{(1)}$ 称为双荷 1-形式对称性群。

\para{离散 theta 角。}$\SO(3)_-$ 理论的作用量包含项
\begin{equation}
    \pi \int \frac{\cP(w_2)}{2} \,,
    \tag{3.193}
\end{equation}
其中 $\cP(w_2)$ 是由 $\SO(3)$ 规范丛的第二 Stiefel--Whitney 类 $w_2$ 取 Pontryagin 平方得到的 $\Z_4$ 值 4 次类。它可以表示为
\begin{equation}
    \cP(w_2) = \tilde{w}_2 \cup \tilde{w}_2 - \delta \tilde{w}_2 \cup_1 \tilde{w}_2 \,,
    \tag{3.194}
\end{equation}
其中 $\tilde{w}_2$ 是一个 $\Z_4$ 值 2-上链，它是 $\Z_2$ 值 2-上闭链 $w_2$ 的提升，而 $\cup_1$ 是一种高阶 cup 积，其定义我们请读者参见 [80] 的附录 B。式 (3.193) 中的 $1/2$ 之所以允许，是因为我们额外假设只研究容许自旋结构的时空上的 4d 规范理论；在这种情况下 $\cP(w_2)$ 是偶类，即它的积分只取 $\{0,2\} \subset \Z_4 = \{0,1,2,3\}$ 中的值。

正如普通 theta 角项是场强的平方，项 (3.193) 是动力学 2 次场 $w_2$ 的平方，而 $w_2$ 可以解释为 $\SO(3)$ 规范理论中的分数场强。出于这个原因，(3.193) 被称为离散 theta 角。

正如上文所讨论的，项 (3.193) 可以理解为施行变换 (3.189) 时产生的对作用量的分数贡献。

% [Tutor Note: p.56 式 (3.194) 是 Pontryagin 平方的上链公式：\tilde{w}_2 是 w2 的 Z4 提升，右端在提升改变时只差一个恰当项，故在 H^4(;Z4) 中给同一个类。这里的 \cup_1 积正是 Steenrod 型高阶 cup 积，(3.194) 就是 P(w2) = w2 ⌣ w2 的整提升修正项。]
